%-----------------------------------
% Source de Base rapport ESEO
%
% <Type de document>	\TypeDocument
% <Portée du document>	\PorteeDocument
%
%-----------------------------------

\documentclass[a4paper,11pt,titlepage]{article}

\usepackage[english,francais]{babel}
\usepackage[T1]{fontenc}
\usepackage[utf8]{inputenc}

\usepackage{xspace, graphicx}
\usepackage{amsmath}
\usepackage{amsfonts}
\usepackage{lastpage}
\usepackage{array}
\usepackage{hyperref}
\usepackage{color}
\usepackage{xcolor}
\usepackage{lmodern}

%------------------------------------

\usepackage{fancyhdr}
\newcommand{\TypeDocument}{Cahier de Tests}
\newcommand{\AuteurPrincipal}{Clément Pabst}
\newcommand{\NomGroupeDocument}{Equipe DDMP}
\newcommand{\IntituleDocument}{Module «~COOL, QL~»}
\newcommand{\PorteeDocument}{A l'attention des enseignants du module «~COOL, QL~» 2017/2018.}
\pagestyle{fancy}

%%%%Ajouts Clément Pabst, Début :

\usepackage[shortlabels]{enumitem}
\setlist[itemize]{label=\textbullet}


\hypersetup{                    % parametrage des hyperliens
    colorlinks=true,                % colorise les liens
    breaklinks=true,                % permet les retours à la ligne pour les liens trop longs
    urlcolor= black,                 % couleur des hyperliens
    linkcolor= darkgray,                % couleur des liens internes aux documents (index, figures, tableaux, equations,...)
    citecolor= black                % couleur des liens vers les references bibliographiques
    pdfauthor={\AuteurPrincipal},     % author
    pdfcreator={\AuteurPrincipal},     % creator
    colorlinks=true,       % false: boxed links; true: colored links
    pdfsubject={Cahier de Test},   % subject of the document
    pdftitle={MonkeyIsland - Equipe DDMP},  %title of the document
    pdfproducer={}
}

%%%%Ajouts Clément Pabst,  Fin  :

\setlength{\hoffset}{-40pt}

\setlength{\topmargin}{-25pt}
\setlength{\headsep}{10pt}

\renewcommand{\headheight}{80pt}

\renewcommand{\headwidth}{450pt}
\setlength{\textwidth}{450pt}
\setlength{\textheight}{604pt}

\renewcommand{\footrulewidth}{0.1mm}

\fancyhf{}
        \fancyhead[LO]{\bf \includegraphics[width=60pt]{img/logoeseo.eps} \textsc{Angers}\\
                                Cycle Ingénieur $3^{\text{éme}}$ année\\
                                \NomGroupeDocument}
         \fancyhead[RO]{\bf {\large \IntituleDocument}\\
				\medskip
				{\large \TypeDocument}\\
				\medskip
				{\small\textit{version du \today}}}
         \fancyfoot[LO]{\sl \PorteeDocument}
         \fancyfoot[RO]{\thepage/\pageref{LastPage}}

\setcounter{tocdepth}{3}

%----------------------------------------
%       DOCUMENT
%----------------------------------------

\begin{document}

%---------------------------------------
\setlength\parindent{0pt}
\sloppy%
\renewcommand{\arraystretch}{1.5}

%---------------------------------------

\vspace{-2cm}%
\begin{center}%
{\Large {\textsc{\bf \IntituleDocument}}}
\vspace{0.4cm}\\
{\large\bf \NomGroupeDocument}

\rule[0.5ex]{0.52\textwidth}{0.1mm}

\vspace{0.2cm}
{\Large\bf{\TypeDocument}}
\vspace{0.2cm}\\
{\Large\bf\textit{\PorteeDocument}}
\end{center}

\vspace{-0.2cm}
\noindent\rule[0.5ex]{\textwidth}{0.1mm}
\textit{Ce document est un cahier de tests dédié au projet «~Monkey Island~». Il est à destination des professeurs en charge de l'évaluation du module de l'année 2017-2018.}\\
\rule[0.5ex]{\textwidth}{0.1mm}

\vspace{1cm}
\begin{tabular}{|p{5cm}|p{9cm}|}
\hline
\textbf{Auteur} & \textbf{Contact} \\\hline
\AuteurPrincipal & clement.pabst@reseau.eseo.fr\\\hline
\end{tabular}

\vspace{0.5cm}
\begin{tabular}{|p{5cm}|p{9cm}|}\hline
\textbf{Relecteurs} & \textbf{Contacts} \\\hline
François de Broch d'Hotelans & francois.debrochdhotelans@reseau.eseo.fr\\\hline
Cailyn Davies & cailyn.davies@reseau.eseo.fr\\\hline
Guillaume Muret & guillaume.muret@reseau.eseo.fr\\\hline
\end{tabular}

%\vspace{1cm}
%\begin{tabular}{|p{1.5cm}|p{2cm}|p{10.1cm}|}\hline
%\textbf{Version} & \textbf{Date} & \textbf{Commentaire} \\\hline
%<x.y> & <date> & Ajout<s> : <sections et/ou sous-sections ajoutées> \newline Suppression<s> : <sections et/ou sous-sections supprimées> \newline Modification<s> : <sections et/ou sous-sections modifiées> \newline Mise<s> à jour : <sections et/ou sous-sections mises à jour> \\
%<...> & <...> & <...> \\
%<1\iere{}> & <date> & <sections présentes> \\\hline
%\end{tabular}
\newpage

%------------------------------- 

\tableofcontents
% alternative pour réduire l'espacement entre les entrées de la table des matières
% (la valeur numérique peut être adaptée au besoin) : 
%{\setlength{\baselineskip}{0.96\baselineskip}\tableofcontents\par}
\newpage

%-------------------------------

%-----------------------------------
% Projet Systèmes Critiques
%
% <Type de document>
% <Portée du document>
% - Introduction
%-----------------------------------

\section{Introduction}
\label{sec:introduction}

\subsection{Objectif}
\label{subsec:introduction:objectif}

L'objectif de ce document est de présenter les tests réalisés dans le cadre du projet Monkey Island. On y vérifiera le bon déroulement des fonctionnalités évoquées dans le plan de test.

\subsection{Portée}
\label{subsec:introduction:portee}

Ce document est à destination des professeurs en charge de l'évaluation du module «~COOL, QL~» de l'année 2017-2018. Il n'a pour vocation qu'une exploitation pédagogique, interne à l'ESEO.




%%-----------------------------------
% Projet Systèmes Critiques
%
% <Type de document>
% <Portée du document>
% - Références
%-----------------------------------

\section{Références}
\label{sec:references}

\subsection{Normes et standards}
\label{subsec:reference:normes}

\begin{tabular}{|l|l|l|}
\hline
\textbf{Norme ou standard} & \textbf{Version} & \textbf{Source} \\
\hline
<nom du document> & <version> & <auteur(s) ou organisme(s)>\\
\hline
\end{tabular} 


\subsection{Documents externes}
\label{subsec:references:externes}


\begin{tabular}{|l|l|l|}
\hline
\textbf{Document} & \textbf{Version} & \textbf{Source} \\
\hline
<nom du document> & <version> & <auteur(s) ou organisme(s)>\\
\hline
\end{tabular} 


\subsection{Documents internes}
\label{subsec:references:internes}


\begin{tabular}{|l|l|l|l|}
\hline
\textbf{Document} & \textbf{Version} & \textbf{Source} \\
\hline
<nom du document> & <version> & <localisation du document>\\
\hline
\end{tabular} 



%\newpage

% \newpage
% \input{Section1}

% \newpage
% \input{Section2}

\newpage
%-----------------------------------
% Projet Systèmes Critiques
%
% <Type de document>
% <Portée du document>
% - <titre section>
%-----------------------------------

\section{Pirates}
\label{sec:pirates}

\subsection{Déplacements}
\label{subsec:piratesDeplacements}

Le but ici va être de valider l'ensemble des déplacements des pirates, incluant leurs interactions avec les autres éléments du jeu (autres personnages, objets et fin de partie).\\
Au cours de ses tests, on peut tester visuellement que les singes respectent le cahier des charges.

\subsubsection{Déplacements de base d'un Pirate}\label{subsubsec:depBasePirate}
\textbf{Résumé :} S'assurer que l'ensemble des déplacements d'un pirate fonctionne. Tous les déplacements sont faits sans détours.\\
\textbf{Précondition :} Pirate placé, entouré de cases terre vides. Une case mer est présente plus loin ainsi qu'un autre pirate qui restera immobile. Le trésor est en place lui aussi, à un endroit connu du testeur. L'ensemble des déplacements nécessaire au scénario est inférieur à l'énergie de départ du pirate.
\begin{table}[!h]
	\begin{center}
		\begin{tabular}{|p{6pt}|c|c|}%p{10cm}
			\hline
			& \textbf{Etape} & \textbf{Résultat(s) Attendu(s)} \\\hline
			1 & Presser successivement H-B-G-D & Déplacement du pirate H-B-G-D  \\\hline
			2 & Presser simultanément HD||DB||BG||GH & Pas de déplacement \\\hline
			3 & Déplacer le pirate sur la case mer & Pas de déplacement \\\hline
			4 & Déplacer le pirate sur l'autre pirate & Pas de déplacement \\\hline
			5 & Déplacer le pirate en dehors de la map & Pas de déplacement \\\hline
			6 & Déplacer le pirate sur le trésor & Victoire et redémarrage de la partie \\\hline
		\end{tabular}
	\end{center}
	\caption{Scénario de déplacement d'un Pirate, obstacle et victoire}
\end{table}

\newpage
\subsubsection{Déplacements \& Interactions d'un Pirate avec d'autres éléments du jeu}
\textbf{Résumé :} S'assurer que la rencontre avec un singe tue le pirate.\\
\textbf{Précondition :} Pirate placé, entouré de cases terre vides. Une bouteille de rhum est placée, visible et pleine, juste à côté du pirate. Un singe chasseur et un singe erratique sont présents sur la carte.\\
\begin{table}[!h]
	\begin{center}
		\begin{tabular}{|p{6pt}|c|c|}
			\hline
			& \textbf{Etape} & \textbf{Résultat(s) Attendu(s)} \\\hline
			1 & Déplacer le pirate sur & Mort du pirate et\\
			 & le singe erratique/chasseur & redémarrage de la partie \\\hline
		\end{tabular}
	\end{center}
	\caption{Scénario de déplacement d'un Pirate, rencontre avec un singe erratique/chasseur}
\end{table}

\textbf{Précondition supplémentaire :} Un autre pirate est présent et se contente d'éviter les singes sans autre interraction.
\begin{table}[!h]
	\begin{center}
		\begin{tabular}{|p{6pt}|c|c|}
			\hline
			& \textbf{Etape} & \textbf{Résultat(s) Attendu(s)} \\\hline
			1 & Déplacer le pirate sur le rhum & Augmentation de l'énergie et\\
			 &  & changement d'état (DRUNK) \\\hline
			2 & Attendre un singe (rester statique) & Mort du pirate \\\hline
			3 & Attente de la réapparition du rhum & Réapparition du rhum, pas de\\
			 &  & changement de l'état(DEAD) du pirate \\\hline
			4 & Déplacer le pirate & Pas de déplacement \\\hline
			5 & Le second pirate fini par s'arrêter & Le singe chasseur tue le pirate \\
			 & & et la partie redémarre. \\\hline
		\end{tabular}
	\end{center}
	\caption{Scénario de déplacement d'un Pirate, rencontre avec un singe et mort sur rhum}
\end{table}

\begin{table}[!h]
	\begin{center}
		\begin{tabular}{|p{6pt}|c|c|}%p{10cm}
			\hline
			& \textbf{Etape} & \textbf{Résultat(s) Attendu(s)} \\\hline
			1 & Déplacer le pirate à côté du trésor & Rien \\
			 & sans autre interraction &  \\\hline
			 & Attendre que le singe chasseur soit sur & Découverte du trésor,\\
			2 & la case du trésor puis déplacer le pirate &   fin de partie\\
			 & sur la case contenant le trésor et le singe &  et mort du pirate\\\hline
		\end{tabular}
	\end{center}
	\caption{Scénario de déplacement d'un Pirate, fin de partie et mort simultanées}
\end{table}

\newpage
\subsection{Energie}
\label{subsec:piratesEnergie}
\textbf{Résumé :} Le but ici va être de valider les conséquences d'un manque d'énergie du pirate ainsi que ses rencontres avec du rhum selon son niveau d'énergie.\\
\textbf{Précondition :} Pirate placé, une carte peuplée uniquement de cases terre vides ainsi que de trois bouteilles de rhum.
\begin{table}[!h]
	\begin{center}
		\begin{tabular}{|p{6pt}|c|c|}%p{10cm}
			\hline
			& \textbf{Etape} & \textbf{Résultat(s) Attendu(s)} \\\hline
			1 & Déplacer le pirate sur une bouteille avec & Déplacements suivants normaux \\
			 & l'énergie<ENERGIE MAX après recharge & \\\hline

			2 & Déplacer le pirate sur une bouteille avec & Déplacements suivants normaux \\
			 & l'énergie=ENERGIE MAX après recharge & \\\hline
			3 & Déplacer le pirate sur une bouteille avec & Déplacements suivants aléatoires jusqu'à une \\
			 & l'énergie>ENERGIE MAX après recharge & baisse d'énergie => énergie=ENERGIE MAX\\\hline
			4 & Déplacer le pirate jusqu'à énergie=0 & Mort du pirate \\\hline
		\end{tabular}
	\end{center}
	\caption{Scénario de déplacement d'un Pirate, ainsi que son rapport à l'énergie}
\end{table}

\newpage
%-----------------------------------
% Projet Systèmes Critiques
%
% <Type de document>
% <Portée du document>
% - <titre section>
%-----------------------------------

\section{Partie}
\label{sec:partie}
Le but ici va être de valider l’ensemble des éléments d'une partie d'un point de vue d'un client.

\subsection{Gestion de la Partie}
\label{subsec:partieGestion}

Tester visuellement lors des scénarios Pirates (cf p.\pageref{subsec:piratesDeplacements}).\\
\begin{itemize}
\item Partie finie -> redémarre si le pirate va sur le trésor,
\item Partie finie si plus de pirate sur la carte,
\item Partie finie -> redémarre si plus de pirate vivant sur la carte.
\end{itemize}

\subsection{Communication}
\label{subsec:partieCommunication}
Les tests sont effectués automatiquement dans JMeter.
\begin{itemize}
\item Inscription du pirate,
\item Déplacements HBGD (4tests),
\item NON-Déplacement de 2 cases HBGD (4tests),
\item NON-Déplacement en diagonale (4tests),
\item NON-Déplacement de 2 cases en diagonale (4tests).
\end{itemize}

\newpage
\input{Analyses}

%\newpage
%%-----------------------------------
% Projet Systèmes Critiques
%
% <Type de document>
% <Portée du document>
% - Définitions
%-----------------------------------

\section{Définitions des termes et acronymes}
\label{sec:definitionsacronymes}

\subsection{Définitions}
\label{subsec:definitionsacronymes:definitions}



\subsection{Acronymes}
\label{subsec:definitionsacronymes:acronymes}






\end{document}

%--- END 

